\documentclass[11pt]{article}
\usepackage[margin=1in]{geometry}
\usepackage{enumitem}
% =============================================================================
% Preambles/header.tex — shared LaTeX/Beamer preamble for this project
%
% Use from a Beamer lecture by prepending:
%     \input{header}
% and compiling with TEXINPUTS=../Preambles:$TEXINPUTS (the /compile-latex
% skill does this automatically).
%
% PALETTE CONTRACT. Color names defined here MUST match the names in
% Quarto/theme-template.scss so Beamer and Quarto renderings use the same
% palette. The scripts/check-palette-sync.sh script enforces this.
%
% To customize: edit the HEX values below AND the matching $variable in
% Quarto/theme-template.scss, then run scripts/check-palette-sync.sh.
% =============================================================================

% -----------------------------------------------------------------------------
% Engine detection + encoding
%
% Our /compile-latex skill uses XeLaTeX, where inputenc is unnecessary and
% can produce encoding warnings. Only load inputenc under pdfTeX.
% -----------------------------------------------------------------------------
\usepackage{iftex}
\ifPDFTeX
  \usepackage[utf8]{inputenc}
\fi

\usepackage{amsmath, amssymb, amsthm}
\usepackage{booktabs}
\usepackage{graphicx}

% -----------------------------------------------------------------------------
% PALETTE — mirrors Quarto/theme-template.scss variable names.
% Load xcolor and define colors BEFORE anything that references them
% (notably hyperref, hypersetup, and the Beamer color assignments).
% -----------------------------------------------------------------------------
\usepackage{xcolor}

\definecolor{primary-blue}{HTML}{012169}      % headings, accents
\definecolor{primary-gold}{HTML}{B9975B}      % emphasis, borders
\definecolor{highlight-yellow}{HTML}{F2A900}  % markers, alerts
\definecolor{light-bg}{HTML}{E8EDF5}          % title / section backgrounds
\definecolor{jet}{HTML}{1A1A1A}               % body text

% Semantic colors (used in diagrams and annotations)
\definecolor{positive}{HTML}{15803D}          % good / observed / identified
\definecolor{negative}{HTML}{B91C1C}          % bad / problematic / bias
\definecolor{neutral}{HTML}{525252}           % reference / context

% Highlight accents (match .hi-* classes in the SCSS)
\definecolor{hi-slate}{HTML}{314F4F}
\definecolor{hi-green}{HTML}{2E8B57}
\definecolor{hi-red}{HTML}{C41E3A}

% Semantic aliases so skill templates and snippets can write
% \textcolor{accent}{...} without hard-coding "primary-blue".
\colorlet{accent}{primary-blue}
\colorlet{accent2}{primary-gold}

% -----------------------------------------------------------------------------
% Hyperref — loaded AFTER xcolor + palette so link colors resolve.
% -----------------------------------------------------------------------------
\usepackage{hyperref}
\hypersetup{colorlinks=true, linkcolor=primary-blue, urlcolor=primary-blue,
            citecolor=primary-blue, pdfborder={0 0 0}}

% -----------------------------------------------------------------------------
% Course code — set once per deck (after \input{header}, before \begin{document})
% via \coursecode{CS401}. Shown in the footer on every slide so a deck's course
% is identifiable without opening the title page. Empty by default.
% -----------------------------------------------------------------------------
\makeatletter
\newcommand{\coursecode}[1]{\gdef\@coursecodetext{#1}}
\gdef\@coursecodetext{}
\makeatother

% -----------------------------------------------------------------------------
% Beamer theme assignments (only apply under Beamer).
%
% We detect Beamer via \@ifclassloaded{beamer}, which is the documented,
% non-internal way. This must be wrapped in \makeatletter/\makeatother
% because \@ifclassloaded contains an @-catcode character.
% -----------------------------------------------------------------------------
\makeatletter
\@ifclassloaded{beamer}{%
  \usefonttheme{professionalfonts}
  \setbeamercolor{structure}{fg=primary-blue}
  \setbeamercolor{frametitle}{fg=primary-blue}
  \setbeamercolor{title}{fg=primary-blue}
  \setbeamercolor{subtitle}{fg=primary-gold}
  \setbeamercolor{author}{fg=primary-blue}
  \setbeamercolor{institute}{fg=neutral}
  \setbeamercolor{date}{fg=primary-gold}
  \setbeamercolor{itemize item}{fg=primary-blue}
  \setbeamercolor{itemize subitem}{fg=primary-gold}
  \setbeamercolor{alerted text}{fg=primary-gold}
  \setbeamercolor{block title}{fg=white, bg=primary-blue}
  \setbeamercolor{block body}{bg=light-bg}
  % Semantic block variants: block=definition/concept (blue), exampleblock=
  % worked example (green/positive), alertblock=key takeaway or caution (gold).
  % Keep to <=2 colored boxes per slide across ALL three variants (INV-7).
  \setbeamercolor{block title example}{fg=white, bg=positive}
  \setbeamercolor{block body example}{bg=light-bg}
  \setbeamercolor{block title alerted}{fg=white, bg=primary-gold}
  \setbeamercolor{block body alerted}{bg=light-bg}

  % Minimal footer: course code (left) + page X / N (right), no navigation clutter
  \setbeamertemplate{navigation symbols}{}
  \setbeamertemplate{footline}{%
    \color{neutral}\footnotesize\hspace{0.7em}\@coursecodetext\hfill
    \insertframenumber/\inserttotalframenumber\hspace{0.7em}\vspace{0.3em}}
}{}
\makeatother

% -----------------------------------------------------------------------------
% TikZ setup — libraries the snippet gallery uses
% -----------------------------------------------------------------------------
\usepackage{tikz}
\usetikzlibrary{arrows.meta, positioning, calc, decorations.pathreplacing,
                fit, shapes.geometric, backgrounds}

% Shared TikZ styles — snippets and hand-written diagrams can pick these up
% via `\begin{tikzpicture}[every node/.style={...}]` or by naming specific
% styles (e.g., `\node[dag-node] (X) at (0,0) {X};`).
\tikzset{
  % Node styles for causal DAGs and similar
  dag-node/.style = {
    draw=primary-blue, fill=light-bg, rounded corners=2pt,
    minimum width=1.2cm, minimum height=0.9cm, align=center, font=\small
  },
  decision-node/.style = {
    draw=primary-gold, fill=white, diamond, aspect=1.8,
    minimum width=1.8cm, minimum height=1.2cm, align=center, font=\small,
    inner sep=1pt
  },
  % Edge styles — observed vs counterfactual
  observed-edge/.style = {-{Latex[length=2.5mm]}, thick, draw=jet},
  counterfactual-edge/.style = {-{Latex[length=2.5mm]}, thick, draw=neutral, dashed},
  confound-edge/.style = {-{Latex[length=2.5mm]}, thick, draw=negative, dashed},
  % Data-point styles for plots
  observed-dot/.style = {fill=primary-blue, draw=primary-blue, circle, inner sep=1.2pt},
  counterfactual-dot/.style = {fill=white, draw=neutral, circle, inner sep=1.2pt}
}

% -----------------------------------------------------------------------------
% Convenience macros
% -----------------------------------------------------------------------------
\newcommand{\muted}[1]{\textcolor{neutral}{#1}}
\newcommand{\key}[1]{\textcolor{primary-gold}{\textbf{#1}}}
\newcommand{\good}[1]{\textcolor{positive}{#1}}
\newcommand{\bad}[1]{\textcolor{negative}{#1}}

% Standout frame macro: full-bleed dark background for section transitions.
% Beamer-only (uses \begin{frame}/\setbeamercolor/\usebeamerfont) -- do not
% call from an article-class file (e.g. Notes/<CODE>/*-notes.tex). \muted,
% \key, \good, \bad, and \coursecode above are plain xcolor/gdef and are
% safe to use from either class; verified when Notes/article-class support
% was added.
% Expand: \transitionslide{Your transition title}
\newcommand{\transitionslide}[1]{%
  \begingroup
  \setbeamercolor{background canvas}{bg=primary-blue}
  \begin{frame}[plain]
    \centering
    \vfill
    {\usebeamerfont{title}\color{white}\Large #1\par}
    \vfill
  \end{frame}
  \endgroup
}

% Section divider frame (Beamer-only), an alternative to \transitionslide with a
% darker two-line style: near-black (jet) background, a small white label line
% above a larger gold title, and a thin gold rule along the bottom edge.
% Expand: \sectiondivider{Section 2}{Pointers}
\newcommand{\sectiondivider}[2]{%
  \begingroup
  \setbeamercolor{background canvas}{bg=jet}
  \begin{frame}[plain]
    \vspace*{3.0cm}
    {\color{white}\ttfamily\large #1\par}
    \vspace{0.5cm}
    {\color{primary-gold}\LARGE #2\par}
    \begin{tikzpicture}[remember picture, overlay]
      \draw[primary-gold, line width=2.5pt]
        ($(current page.south west)+(0,0.1)$) -- ($(current page.south east)+(0,0.1)$);
    \end{tikzpicture}
  \end{frame}
  \endgroup
}

% -----------------------------------------------------------------------------
% Channel watermark — "Concepts That Click" (YouTube).
%
% Article-class files (CompetitiveExam Q/A sets, Notes, Assignments) get a
% light diagonal draft watermark on every page plus a centered footer naming
% the channel. Beamer decks get a subtle TikZ background watermark instead
% (the footline already carries the course code). Centralized here so every
% MCQ PDF is branded without per-file edits.
% -----------------------------------------------------------------------------
\makeatletter
\@ifclassloaded{article}{%
  \usepackage{draftwatermark}
  \SetWatermarkText{Concepts That Click}
  \SetWatermarkScale{1}
  \SetWatermarkFontSize{2cm}
  \SetWatermarkLightness{0.86}
  \SetWatermarkAngle{45}
  \usepackage{fancyhdr}
  \pagestyle{fancy}
  \fancyhf{}
  \fancyfoot[C]{\footnotesize\textcolor{neutral}{Concepts That Click~|~YouTube: \href{https://www.youtube.com/@concepts.thatclick}{@concepts.thatclick}}}
  \renewcommand{\headrulewidth}{0pt}
  \renewcommand{\footrulewidth}{0pt}
  \fancypagestyle{plain}{%
    \fancyhf{}%
    \fancyfoot[C]{\footnotesize\textcolor{neutral}{Concepts That Click~|~YouTube: \href{https://www.youtube.com/@concepts.thatclick}{@concepts.thatclick}}}%
    \renewcommand{\headrulewidth}{0pt}%
    \renewcommand{\footrulewidth}{0pt}%
  }
}{}
\@ifclassloaded{beamer}{%
  \setbeamertemplate{background}{%
    \begin{tikzpicture}[remember picture, overlay]
      \node[rotate=45, opacity=0.055, align=center]
        at (current page.center)
        {\Huge\textcolor{neutral}{Concepts That Click}};
    \end{tikzpicture}%
  }
}{}
\makeatother 
\coursecode{JKSSB}
\title{Lesson 08 Practice: RAM, ROM \& Cache}
\author{JKSSB Computer Awareness}
\date{\today}
\begin{document}
\maketitle
\noindent\textbf{Provenance:} All 20 items are original JKSSB-pattern
questions written for this lesson. None is a real PYQ; each is labelled
\textit{[Original]}.

\begin{enumerate}[label=\textbf{Q\arabic*.}, leftmargin=*]
\item \textit{[Original]} Which type of memory loses its contents when the
  computer is switched off?
  \begin{enumerate}[label=\Alph*.]
    \item ROM
    \item Hard disk
    \item RAM
    \item Flash drive
  \end{enumerate}
\item \textit{[Original]} RAM stands for:
  \begin{enumerate}[label=\Alph*.]
    \item Random Access Memory
    \item Read Access Memory
    \item Rapid Access Memory
    \item Read-Only Access Memory
  \end{enumerate}
\item \textit{[Original]} ROM stands for:
  \begin{enumerate}[label=\Alph*.]
    \item Read-Only Memory
    \item Random-Only Memory
    \item Read-Output Memory
    \item Rapid-Only Memory
  \end{enumerate}
\item \textit{[Original]} Which statement is \textbf{NOT} correct?
  \begin{enumerate}[label=\Alph*.]
    \item RAM is volatile.
    \item ROM is non-volatile.
    \item RAM retains its contents after power-off.
    \item ROM is read-only in normal operation.
  \end{enumerate}
\item \textit{[Original]} In the office desktop (EX-01), the main memory that
  holds the running programs and is cleared when the machine is switched off
  is:
  \begin{enumerate}[label=\Alph*.]
    \item ROM
    \item RAM
    \item L1 cache
    \item the SSD
  \end{enumerate}
\item \textit{[Original]} Which of the following is volatile memory?
  \begin{enumerate}[label=\Alph*.]
    \item DRAM
    \item ROM
    \item PROM
    \item EPROM
  \end{enumerate}
\item \textit{[Original]} Which memory is faster and does \textbf{not}
  require periodic refreshing?
  \begin{enumerate}[label=\Alph*.]
    \item DRAM
    \item ROM
    \item Hard disk
    \item SRAM
  \end{enumerate}
\item \textit{[Original]} DRAM differs from SRAM mainly because DRAM:
  \begin{enumerate}[label=\Alph*.]
    \item is faster and needs no refresh
    \item is used for cache memory
    \item is non-volatile
    \item is denser and cheaper but needs periodic refresh
  \end{enumerate}
\item \textit{[Original]} SRAM is commonly used for:
  \begin{enumerate}[label=\Alph*.]
    \item main memory
    \item cache memory
    \item secondary storage
    \item firmware
  \end{enumerate}
\item \textit{[Original]} DRAM is commonly used for:
  \begin{enumerate}[label=\Alph*.]
    \item main memory
    \item cache memory
    \item permanent storage of the start-up code
    \item CPU registers
  \end{enumerate}
\item \textit{[Original]} Which abbreviation--full form pairing is correct?
  \begin{enumerate}[label=\Alph*.]
    \item DRAM --- Direct Random Access Memory
    \item EEPROM --- Electrically Erasable Programmable Read-Only Memory
    \item EPROM --- Electrically Programmable Read-Only Memory
    \item SRAM --- Serial Random Access Memory
  \end{enumerate}
\item \textit{[Original]} Which ROM type can be programmed once by the user
  and cannot be erased?
  \begin{enumerate}[label=\Alph*.]
    \item EPROM
    \item EEPROM
    \item DRAM
    \item PROM
  \end{enumerate}
\item \textit{[Original]} Which ROM type is erased by exposing it to
  ultraviolet light?
  \begin{enumerate}[label=\Alph*.]
    \item PROM
    \item EPROM
    \item EEPROM
    \item Mask ROM
  \end{enumerate}
\item \textit{[Original]} Which memory can be erased and rewritten
  electrically, without removing it from the circuit?
  \begin{enumerate}[label=\Alph*.]
    \item PROM
    \item EPROM
    \item EEPROM
    \item SRAM
  \end{enumerate}
\item \textit{[Original]} Flash memory is best described as:
  \begin{enumerate}[label=\Alph*.]
    \item a volatile form of DRAM
    \item an EEPROM-derived non-volatile memory erased in blocks
    \item the mask-programmed type of ROM
    \item a type of CPU register
  \end{enumerate}
\item \textit{[Original]} Evaluate the statements:
  \begin{enumerate}[label=\Roman*.]
    \item RAM is volatile, while ROM is non-volatile.
    \item ROM is read-only in normal operation, though special programming
      is possible.
    \item Cache memory is slower than main memory.
  \end{enumerate}
  \begin{enumerate}[label=\Alph*.]
    \item I and II only
    \item II and III only
    \item I and III only
    \item I, II, and III
  \end{enumerate}
\item \textit{[Original]} Which cache level is the smallest, fastest, and
  closest to the CPU core?
  \begin{enumerate}[label=\Alph*.]
    \item L1
    \item L2
    \item L3
    \item Main memory
  \end{enumerate}
\item \textit{[Original]} Which cache level is the largest and is typically
  shared across the processor's cores?
  \begin{enumerate}[label=\Alph*.]
    \item L1
    \item L2
    \item ROM
    \item L3
  \end{enumerate}
\item \textit{[Original]} In cache operation, a ``miss'' occurs when:
  \begin{enumerate}[label=\Alph*.]
    \item the required data is found in the cache
    \item data is written to ROM
    \item the required data is not found in the cache
    \item the cache is powered off
  \end{enumerate}
\item \textit{[Original]} A cache handles 100 memory accesses, of which 90
  are hits. The hit ratio is:
  \begin{enumerate}[label=\Alph*.]
    \item 0.10
    \item 1.00
    \item 9.0
    \item 0.90
  \end{enumerate}
\end{enumerate}
\end{document}
